Para resumir cuánto se parece el orden de los algoritmos entre paper y réplica usamos la tau de Kendall. Compara las dos clasificaciones pareja por pareja: más uno es idéntico, cero es sin relación, menos uno es al revés.

Sale siempre positiva: el orden siempre se parece. El desacuerdo de los cinco primeros escenarios viene casi todo de Folly, el caso que ya hemos visto. Y el valor bajo del último es el empate de la RAM: cuando cuatro van igualados, el orden entre ellos es casi azar.
